\section{Introduction}

Quantum low-density parity-check (qLDPC) codes promise fault-tolerant memory at a fraction of the
surface code's qubit overhead \citep{bravyi2024gross}, and belief-propagation decoders --- in
particular relay-BP \citep{muller2025relaybp} --- have been shown, in software, to reach the logical
error rates that make such codes useful without an ordered-statistics tail. Whether that translates
into a decoder that keeps up with a quantum processor is a hardware question: the decode of one
syndrome window must complete in about a microsecond for superconducting platforms, and in the
tens to hundreds of microseconds for ions and neutral atoms, on a budget of area and power that a
control system can absorb.

This paper documents one end-to-end answer, built in the open and measured at every step. The
contributions, in the order the work was done:

\begin{enumerate}
  \item A fixed-point relay-BP golden model for the \gross{} gross code that is the RTL specification:
        an 8-bit Q5.3 message word and a $6\times10$ schedule, each chosen by a sweep against the
        floating-point decoder on circuit-level noise (\S\ref{sec:decoder}).
  \item A ladder of synthesizable cores from a sequential FSM to a $K$-banked, constant-geometry
        design whose gather/scatter is a fixed permutation network and whose message storage is a
        latch register file, with a co-simulation harness that drives every variant --- and the
        routed ASIC netlist --- with the same golden vectors (\S\ref{sec:architecture}).
  \item Silicon results on a \$300 Kria KV260: \SI{15.64}{\micro\second} worst-case,
        \SI{0.85}{\micro\second} median, bit-exact over $3\times10^6$ shots, and the finding that the
        convergence flag heralds essentially every logical error (\S\ref{sec:silicon}).
  \item The full-parallel geometry on a rented Virtex UltraScale+ and on ASAP7 through OpenROAD:
        543 cycles, \SI{0.88}{\micro\second} at 7\,nm predictive, with the first power figure derived
        from gate-level switching activity rather than default toggle rates (\S\ref{sec:scaling}).
  \item The negative results that shaped the design and that a tape-out would have to face: no OSD
        tail, no feasible \texttt{NGROUP} setting, a streaming core that does not fit, and a
        latch-clock hold failure that repair cannot buffer away (\S\ref{sec:limitations}).
\end{enumerate}

Every number in this paper is copied from a dated report in the public repository
\citep{alephrepo}; Appendix~\ref{app:prov} maps each to its source file. Nothing measured on ASAP7
is called ``the chip'': it is a 7\,nm \emph{predictive} kit, and the target node is not.
